\heading{数学物理方法（下）-第13次作业}


\begin{question}
求下列常微分方程边值问题相应的 Green 函数：
\[
\frac d{dx}\left[(1-x^2)y'\right]=-\phi(x),\qquad
 y(-1),y(1)\text{ 有界}.
\]
\begin{solution}
算子有常数零模，所以需满足相容条件
\[
\int_{-1}^1\phi(x)\,dx=0.
\]
在零均值意义下令
\[
L_xG(x,\xi)=-\delta(x-\xi)+\frac12.
\]
由跳跃条件
\[
\left[(1-x^2)G_x\right]_{\xi-}^{\xi+}=-1
\]
及端点有界性可得 Green 函数。常数规范任意；取零均值规范时，解写为
\[
y(x)=\int_{-1}^1G(x,\xi)\phi(\xi)\,d\xi .
\]
\end{solution}
\end{question}

\begin{question}
求一维和二维无界空间 Helmholtz 方程散射问题的 Green 函数。时间因子取 \(e^{-i\omega t}\)。
\begin{solution}
出射 Green 函数满足
\[
(\nabla^2+k^2)G=-\delta.
\]
一维中在 \(x\ne x'\) 处有
\[
G''+k^2G=0.
\]
出射条件要求 \(x>x'\) 时取 \(e^{ik(x-x')}\)，\(x<x'\) 时取 \(e^{ik(x'-x)}\)，所以设
\[
G=C e^{ik|x-x'|}.
\]
在 \(x=x'\) 附近积分方程，得到导数跳跃
\[
G_x(x'+0,x')-G_x(x'-0,x')=-1.
\]
代入上式得 \(2ikC=-1\)，因此
\[
G_1(x,x')=\frac{i}{2k}e^{ik|x-x'|}.
\]
二维中因平移和旋转不变性，\(G\) 只依赖 \(R=|\bm r-\bm r'|\)。在 \(R>0\) 处径向方程为
\[
G_{RR}+\frac1R G_R+k^2G=0,
\]
其出射解为 \(H_0^{(1)}(kR)\)。由小量展开
\[
H_0^{(1)}(kR)\sim 1+\frac{2i}{\pi}\ln R
\]
与二维基本解的对数奇性匹配，得
\[
G_2(\bm r,\bm r')=\frac{i}{4}H_0^{(1)}(k|\bm r-\bm r'|).
\]
\end{solution}
\end{question}

\begin{question}
求二维 Poisson 方程
\[
\nabla^2G(x,y;x',y')=-4\pi\delta(x-x')\delta(y-y')
\]
的通解。
\begin{solution}
由
\[
\nabla^2\ln |\bm r-\bm r'|=2\pi\delta(\bm r-\bm r')
\]
知一个特解为
\[
G_0=-2\ln |\bm r-\bm r'|.
\]
故通解为
\[
G=-2\ln |\bm r-\bm r'|+H(\bm r,\bm r'),
\]
其中 \(H\) 关于 \(\bm r\) 调和，由边界条件确定。
\end{solution}
\end{question}

\begin{question}
求无界空间二维波动问题
\[
\begin{cases}
u_{tt}-a^2\nabla^2u=f(x,y,t),\qquad -\infty<x,y<\infty,
\quad t>0,\\
u|_{t=0}=\psi(x,y),\qquad u_t|_{t=0}=\xi(x,y)
\end{cases}
\]
相应的 Green 函数。
\begin{solution}
迟滞 Green 函数为
\[
G(\bm r,t;\bm r',t')=
\frac{H(t-t'-R/a)}{2\pi a^2\sqrt{(t-t')^2-R^2/a^2}},
\qquad R=|\bm r-\bm r'|.
\]
这个式子可由空间 Fourier 变换推出。对
\[
G_{tt}-a^2\nabla^2G=\delta(\bm r-\bm r')\delta(t-t')
\]
作空间变换，得到每个波数的受迫振子
\[
\hat G_{tt}+a^2|\bm k|^2\hat G=e^{-i\bm k\cdot\bm r'}\delta(t-t').
\]
迟滞条件给出
\[
\hat G=\frac{\sin(a|\bm k|(t-t'))}{a|\bm k|}H(t-t')e^{-i\bm k\cdot\bm r'}.
\]
反变换后得到上面的二维迟滞核。因二维波动有尾项，核在圆锥内部 \(t-t'>R/a\) 也不为零。
\end{solution}
\end{question}

\begin{question}
求三维无界空间热传导问题的 Green 函数。
\begin{solution}
热方程
\[
G_t-\kappa\nabla^2G=\delta(\bm r-\bm r')\delta(t-t'),\qquad G=0\quad(t<t')
\]
的基本解为
\[
G=\frac{H(t-t')}{[4\pi\kappa(t-t')]^{3/2}}
\exp\left[-\frac{|\bm r-\bm r'|^2}{4\kappa(t-t')}\right].
\]
\end{solution}
\end{question}

\begin{question}
用 Green 函数方法解三维无界空间热传导问题
\[
\begin{cases}
u_t-\kappa\nabla^2u=f(x,y,z,t),\\
u|_{t=0}=\phi(x,y,z).
\end{cases}
\]
\begin{solution}
由三维热核得
\[
u(\bm r,t)=\int_{\mathbb R^3}G(\bm r,t;\bm r',0)\phi(\bm r')\,d\bm r'
+\int_0^t\int_{\mathbb R^3}G(\bm r,t;\bm r',\tau)f(\bm r',\tau)\,d\bm r' d\tau.
\]
其中 \(G\) 为上一题热核。
\end{solution}
\end{question}

\begin{question}
求泛函
\[
J[u]=\iiint_{x^2+y^2+z^2<a^2}\left[(\nabla u)^2-k^2u^2\right]dxdydz
\]
在边界条件 \(u|_{x^2+y^2+z^2=a^2}=z\) 下的极值函数，其中 \(ka\ne\tan ka\)。
\begin{solution}
Euler 方程为
\[
\nabla^2u+k^2u=0.
\]
边界值 \(z=r\cos\theta\) 只含 \(l=1\) 模态，故令
\[
u(r,\theta)=Cj_1(kr)\cos\theta.
\]
由 \(u(a,\theta)=a\cos\theta\) 得
\[
C=\frac{a}{j_1(ka)}.
\]
因此
\[
u(r,\theta)=a\frac{j_1(kr)}{j_1(ka)}\cos\theta.
\]
条件 \(ka\ne\tan ka\) 保证 \(j_1(ka)\ne0\)。
\end{solution}
\end{question}

\begin{question}
写出常微分方程定解问题
\[
\begin{cases}
y''-2xy'+2ny=0,\\
y(0)=0,\qquad y(1)=1
\end{cases}
\]
所对应的泛函极值问题。
\begin{solution}
乘以积分因子 \(e^{-x^2}\)，方程化为
\[
\left(e^{-x^2}y'\right)'+2ne^{-x^2}y=0.
\]
对应泛函可取
\[
J[y]=\int_0^1 e^{-x^2}\left[(y')^2-2ny^2\right]dx,
\]
边界条件为 \(y(0)=0,\ y(1)=1\)。
\end{solution}
\end{question}

\begin{question}
求泛函
\[
J[y]=\int_0^1\left[\left(x\frac{dy}{dx}\right)^2+2y^2\right]dx
\]
在边界条件 \(y(0)\) 有界、\(y(1)=0\) 及约束
\[
J_1[y]=\int_0^1y^2(x)x^2\,dx=1
\]
下的最小值。
\begin{solution}
引入 Lagrange 乘子 \(\lambda\)，Euler 方程为
\[
-(x^2y')'+2y=\lambda x^2y,
\]
即
\[
x^2y''+2xy'+(\lambda x^2-2)y=0.
\]
这是球 Bessel 型方程，满足有界性和 \(y(1)=0\) 的基模给出最小值；最小值等于对应的最小本征值 \(\lambda_1\)。
\end{solution}
\end{question}

\begin{question}
写出算符
\[
\hat L^2S(\theta,\phi)=-\left\{\frac1{\sin\theta}\frac\partial{\partial\theta}\left[\sin\theta\frac{\partial S}{\partial\theta}\right]+\frac1{\sin^2\theta}\frac{\partial^2S}{\partial\phi^2}\right\}
\]
在球面周期和端点有界条件下的本征值问题，并写出其本征值和本征函数。
\begin{solution}
本征值问题为
\[
\hat L^2S=\lambda S,
\]
其中 \(S(0,\phi),S(\pi,\phi)\) 有界，且关于 \(\phi\) 以 \(2\pi\) 为周期。分离变量得球谐函数
\[
S(\theta,\phi)=Y_l^m(\theta,\phi),
\]
本征值为
\[
\lambda=l(l+1),\qquad l=0,1,2,\ldots,
\quad m=-l,-l+1,\ldots,l.
\]
\end{solution}
\end{question}

